%Dokumentklasse

%draft als optionohne bilder für bessere performance
%\documentclass[a4paper,12pt,]{scrreprt}

%normal mit Bildern
\documentclass[
a4paper,
12pt,
draft=True]
{scrartcl}

%angepasster \today Command
\newcommand{\leadingzero}[1]{\ifnum #1<10 0\the#1\else\the#1\fi}
\newcommand{\todayDE}{\leadingzero{\day}.\leadingzero{\month}.\the\year}

%Section als Chapter
\RedeclareSectionCommand[%
%beforeskip = -1sp plus -1sp minus -1sp,% kleinster negativer Wert, um den Absatzeinzug nach der Überschrift zu verhindern.
afterskip = 1.5 \baselineskip plus -1sp minus 1sp,
font = \Huge,
]{section}

\usepackage[left= 3cm,right = 3cm, bottom = 3cm,top = 3cm]{geometry}
%\usepackage[onehalfspacing]{setspace}

% ============= Packages =============
% Dokumentinformationen
\usepackage[
pdftitle={Praktikum - Umwelttechnik},
pdfsubject={},
pdfauthor={Roman-Luca Zank},
pdfkeywords={},	
%Links nicht einrahmen
hidelinks
]{hyperref}

%nur Text zum prüfen des Umfangs

% Standard Packages
%\usepackage[bottom]{footmisc}
\usepackage[utf8]{inputenc}
\usepackage[ngerman]{babel}

\usepackage[T1]{fontenc}
%\usepackage{helvet}

%\renewcommand{\familydefault}{\sfdefault}

\usepackage{colortbl}
\usepackage{graphicx}
\graphicspath{{img/}}
\usepackage{mhchem}
\usepackage{fancyhdr}
\usepackage{lmodern}
\usepackage[table]{xcolor}
\usepackage{placeins}
\usepackage{booktabs}
\usepackage{caption}
\usepackage[list=true]{subcaption}
\usepackage{longtable}
\usepackage{tikz}
\usepackage{pgfplots}
\pgfplotsset{/pgf/number format/use comma}
\pgfplotsset{grid style={white!90!black}}
\usepackage{lastpage}
%\usepackage{ulem}
\usepackage{mathtools}
\usepackage{adjustbox}
\usetikzlibrary{patterns}
\usepackage{pdfpages}
\usepackage[shortlabels]{enumitem}
\usepackage{chemfig}

%Einheitenpackage
\usepackage{siunitx}  
\sisetup{	locale = DE, 
	per-mode=fraction,
	inter-unit-product=\ensuremath{\cdot},
	detect-weight = true,
	quotient-mode=fraction
}
%neue Einheiten definieren
\DeclareSIUnit\xyz{xyz}	
\DeclareSIUnit\rpm{rpm}	
\DeclareSIUnit\mws{mWS}	
\DeclareSIUnit\degrees{^\circ}	

%Automatisch cdot statt *
\DeclareMathSymbol{*}{\mathbin}{symbols}{"01}


%Tabelle
\usepackage{tabularx}
\usepackage{tabulary}

%nur letzte Zeile der Gleichung nummerieren
\makeatletter
\def\Let@{\def\\{\notag\math@cr}}
\makeatother

% zusätzliche Schriftzeichen der American Mathematical Society
\usepackage{amsfonts}
\usepackage{amsmath}

%Abkürzungsverzeichnis
\usepackage{acronym}

%kein Abstand bei neuem Kapitel vom Seitenanfang
%\vspace*{2.3\baselineskip} = ORIGINAL
%\renewcommand*{\chapterheadstartvskip}{\vspace*{.0\baselineskip}}

%nicht einrücken nach Absatz
\setlength{\parindent}{0pt}


\urlstyle{same}

% ============= Kopf- und Fußzeile =============
\pagestyle{fancy}
%
\lhead{}
\chead{}
\rhead{}%\slshape }%\leftmark}
%%
\lfoot{}
\cfoot{}
\rfoot[{\thepage\ of \pageref*{LastPage}}]{Seite \thepage\ von \pageref*{LastPage}}
%%
\renewcommand{\headrulewidth}{0pt}
\renewcommand{\footrulewidth}{0pt}
%\renewcommand{\chapterpagestyle}{fancy}

%Fußnotelinie
%\let\footnoterule

%Fußnote mit Klammer
\renewcommand*{\thefootnote}{(\arabic{footnote})}

%Abb. statt Abbildung
\addto\captionsngerman{%
	\renewcommand{\figurename}{Abb.}%
	\renewcommand{\tablename}{Tab.}%
}

% ============= Package Einstellungen & Sonstiges ============= 
%Besondere Trennungen
%\hyphenation{De-zi-mal-tren-nung}
\usepackage[none]{hyphenat}
\hyphenpenalty=5000
\tolerance=5000
\providecommand\phantomsection{}

\usepackage{mathtools}


% ============= Dokumentbeginn =============

\begin{document}
%Seiten ohne Kopf- und Fußzeile sowie Seitenzahl
\pagestyle{empty}

%\begin{center}
\begin{tabular}{p{\textwidth}}


\begin{center}
\includegraphics[scale=0.75]{logos.jpg}\\
\end{center}


\\

\begin{center}
\LARGE{\textsc{
Protokoll \\
Organische Chemie III\\
}}
\end{center}

\\

%\begin{center}
%\large{Fakultät für Muster und Beispiele \\
%der Hochschule Musterhausen \\}
%\end{center}
%
%\\

\begin{center}
\textbf{\Large{Darstellung des Azofarbstoffes Orange II}
\end{center}


\\
%\begin{center}
%zur Erlangung des akademischen Grades\\
%Bachelor of Engineering
%\end{center}


%\begin{center}
%vorgelegt von
%\end{center}

\begin{center}
\Large{\textbf{Teilnehmer:}} \\ 
\end{center}
\begin{center}
\large{Roman-Luca Zank} \\
\end{center}


\\

\begin{center}
\begin{tabular}{lll}
\large{\textbf{Protokollführer:}} & & \large{Roman-Luca Zank}\\
&&\\
\large{\textbf{Datum der Versuchsdurchführung:}}&& \large{16.12.2019}\\
&&\\
\large{\textbf{Abgabedatum:}}&& \large{\todayDE}
\end{tabular}
\end{center}

\\ \\ \\ \\ \\ 
\large{Merseburg den \todayDE}

\end{tabular}
\end{center}


%\include{14_danksagungen}

%\include{15_zusammenfassung}

% Beendet eine Seite und erzwingt auf den nachfolgenden Seiten die Ausgabe aller Gleitobjekte (z.B. Abbildungen), die bislang definiert, aber noch nicht ausgegeben wurden. Dieser Befehl fügt, falls nötig, eine leere Seite ein, sodaß die nächste Seite nach den Gleitobjekten eine ungerade Seitennummer hat. 
\cleardoubleoddpage

% Pagestyle für Titelblatt leer
\pagestyle{empty}

%Seite zählen ab
\setcounter{page}{0}

%Titelblatt
\begin{center}
\begin{tabular}{p{\textwidth}}


\begin{center}
\includegraphics[scale=0.75]{logos.jpg}\\
\end{center}


\\

\begin{center}
\LARGE{\textsc{
Protokoll \\
Organische Chemie III\\
}}
\end{center}

\\

%\begin{center}
%\large{Fakultät für Muster und Beispiele \\
%der Hochschule Musterhausen \\}
%\end{center}
%
%\\

\begin{center}
\textbf{\Large{Darstellung des Azofarbstoffes Orange II}
\end{center}


\\
%\begin{center}
%zur Erlangung des akademischen Grades\\
%Bachelor of Engineering
%\end{center}


%\begin{center}
%vorgelegt von
%\end{center}

\begin{center}
\Large{\textbf{Teilnehmer:}} \\ 
\end{center}
\begin{center}
\large{Roman-Luca Zank} \\
\end{center}


\\

\begin{center}
\begin{tabular}{lll}
\large{\textbf{Protokollführer:}} & & \large{Roman-Luca Zank}\\
&&\\
\large{\textbf{Datum der Versuchsdurchführung:}}&& \large{16.12.2019}\\
&&\\
\large{\textbf{Abgabedatum:}}&& \large{\todayDE}
\end{tabular}
\end{center}

\\ \\ \\ \\ \\ 
\large{Merseburg den \todayDE}

\end{tabular}
\end{center}
 %Protokolle
%\begin{center}
\begin{tabular}{p{\textwidth}}


\begin{center}
\includegraphics[scale=0.75]{img/logos.jpg}\\
\end{center}


\\

\begin{center}
\LARGE{\textsc{
Recherche \\
Rückgewinnung von Ammoniak aus Industrieabwässern\\
}}
\end{center}

%\begin{center}
%\large{Fakultät für Muster und Beispiele \\
%der Hochschule Musterhausen \\}
%\end{center}
%
%\\
 \\
 
\begin{center}
\textbf{\Large{Seminararbeit in Medienrecherche}}
\end{center}

\begin{center}
	\large{im WiSe 2019}
\end{center}
 \\
%\begin{center}
%zur Erlangung des akademischen Grades\\
%Bachelor of Engineering
%\end{center}


\begin{center}
\large{vorgelegt von}
\end{center}
\\


\begin{center}
\Large{\textbf{Roman-Luca Zank}} \\
\end{center}

\begin{center}
3. Semester \\
Chemie- und Umwelttechnik \\
\end{center}


\begin{center}
\begin{tabular}{lll}
	\textbf{E-Mail:} & & romanzank@mail.de\\
	\textbf{Matrikelnummer:} & &25240\\
	\textbf{Adresse:} & &Platz der Bausoldaten 2, Zimmer 224\\
	\textbf{Ort:} & &06217 Merseburg\\
	&& \\
	\textbf{Prüfer:} & & Dr. Frank  Baumann\\
\end{tabular}
\end{center}

\\ \\ \\ \\ \\
\large{Merseburg, \today}

\end{tabular}
\end{center}
 %Seminar-/Abschlussarbeit

% Pagestyle für Rest des Dokuments
\pagestyle{fancy}

%Inhaltsverzeichnis
%\tableofcontents
%\thispagestyle{empty}
%\newpage

%Inhalt
%
%Verzeichnis aller Bilder
\label{sec:bilder}
\listoffigures
\addcontentsline{toc}{chapter}{Abbildungsverzeichnis}
\thispagestyle{empty}

%Verzeichnis aller Tabellen
\label{sec:tabellen}
\listoftables
\addcontentsline{toc}{chapter}{Tabellenverzeichnis}
\thispagestyle{empty}



%%Abkürzungsverzeichnis
\addsec{Nomenklatur}
\label{sec:abkurzung}


\begin{acronym}[AAA]
			\aacro{kompl}{$P$}{Kompressorleistung (Messgerät)}
			\aacro{kompldia}{$P_{\text{Dia}}$}{Kompressorleistung (Diagramm)}
			\aacro{mKM}{$\dot{m}_{\text{KM}}$}{Massenstrom des Kältemittels}
			\aacro{mKW}{$\dot{m}_{\text{KW}}$}{Massenstrom des Kühlwassers}
			\aacro{dpV}{$\Delta p_1$}{gemessener Verdampferdruck}
			\aacro{pV}{$p_1$}{absoluter Verdampferdruck}
			\aacro{dpK}{$\Delta p_2$}{gemessener Kondensatordruck}
			\aacro{pK}{$p_2$}{absoluter Kondensatordruck}
			\aacro{t1}{$t_1$}{Verdampferaustrittstemperatur}
			\aacro{t2}{$t_2$}{Kondenstoreintrittstemperatur}
			\aacro{t3}{$t_3$}{Kondenstoraustrittstemperatur}
			\aacro{t4}{$t_4$}{Verdampfereintrittstemperatur}	\aacro{t5}{$t_5$}{Eintrittstemperatur Kühlwasser (Kompressor)}
	\aacro{t6}{$t_6$}{Austrittstemperatur Kühlwasser (Kompressor)}
	\aacro{t7}{$t_7$}{Austrittstemperatur Kühlwasser (Kondensator)}
	\aacro{t7}{$t_8$}{Austrittstemperatur Kühlwasser (Verdampfer)}
	\aacro{h}{$h_i$}{Kältemittelenthalpien am Punkt $i$}
	\aacro{tü}{$t_{\text{Ü}}$}{Überhitzungstemperatur}
	\aacro{tkond}{$t_{\text{Kond}}$}{Kondensationstemperatur bei $p_2$}
	\aacro{vol}{$\dot{V}_1$}{Volumenstrom durch den Kompressor}
	\aacro{t9}{$t_{9}$}{Eintrittstemperatur der Kühlluft}
	\aacro{10}{$t_{10}$}{Austrittstemperatur der Kühlluft}
	\aacro{vl}{$v_L$}{Luftgeschwindigkeit am luftgekoppelten Verdampfer}
	\aacro{pu}{$p_u$}{Umgebungsdruck}
	\aacro{A}{$A$}{Fläche des luftgekoppelten Verdampfers}
	\aacro{qwv}{$\dot{Q}_{W,\text{Verd}}$}{Wärme im Verdampfer ausgehend vom Wasser}
	\aacro{q41}{$\dot{Q}_{4-1}$}{Wärme im Verdampfer an das Kältemittel}
	\aacro{q23}{$\dot{Q}_{2-3}$}{Wärme im Kondensator ausgehend vom Kältemittel}
	\aacro{qwk}{$\dot{Q}_{W,\text{kond}}$}{Wärme im Kondensator an das Wasser}
	\aacro{q12}{$\dot{Q}_{1-2}$}{Wärme im Kompressor ausgehend vom Kältemittel}
	\aacro{qw56}{$\dot{Q}_{W, 5-6}$}{Wärme im Kompressor an das Wasser}
	\aacro{qwges}{$\dot{Q}_{W, 5-7}$}{Gesamtwärmeübertragung an das Wasser}
	\aacro{qü}{$\dot{Q}_{Ü}$}{Überhitzungswärme}
	\aacro{qk}{$\dot{Q}_{K}$}{Kondensationswärme}
	\aacro{qk}{$\dot{Q}_{H}$}{Heizwärme}
	\aacro{lz}{$\varepsilon$}{Leistungsziffer}
	\aacro{fracp}{$\pi$}{Druckverhältnis}
	\aacro{dt}{$\Delta T$}{Temperaturhub}
\end{acronym}

%\printnoidxglossary[style=listdotted]

\section{Einleitung und Versuchsziel}
\label{sec:aufgabenstellung}
%In der Aufgabenstellung wird (in eigenen Worten und ganzen Sätzen) formuliert, was das Ziel des 
%Versuches ist.  
%[Beachten Sie die eigentliche Aufgabenstellung in den Versuchsanleitungen sowie die Hinweise zur Auswertung!] 

Im folgenden Versuch wird das Strömungsverhalten variierender Wasservolumenströme durch unterschiedliche Rohrleitungen und Rohrleitungseinbauten untersucht. Dabei ist vor allem die nicht-lineare Abhängigkeit des Druckverlustes in der Rohrleitung und der Geschwindigkeit des Fluides darzustellen. Die sich daraus ergebenden Rohrreibungskoeffizienten $\lambda$ sind in Abhängigkeit von der Reynoldszahl Re zu bestimmen, ebenso wie der Einfluss von Messfehlern auf die Ergebnisse. 

%\section{Physikalische Hintergründe}
\label{sec:physik}

Als physikalische Hintergründe sind die folgenden Gleichungen dargestellt. Diese beschränken sich im Wesentlichen auf die \textsc{Bernoulli}-Gleichung mit deren Ableitungen, die \textsc{Reynold}szahl, sowie dem K$_V$-Wert.\\

\textsc{Bernoulli}-Gleichung:
\begin{flalign}
	p_1+z_1*g*\rho +\frac{1}{2}*\rho*(v_1)^2&= p_2+z_2*g*\rho+\frac{1}{2}*\rho*(v_2)^2
\end{flalign}

\textsc{Bernoulli}-Gleichung mit Druckverlust $\Delta p_v$:
\begin{flalign}
p_1+z_1*g*\rho +\frac{1}{2}*\rho*(v_1)^2&= p_2+z_2*g*\rho+\frac{1}{2}*\rho*(v_2)^2\boldsymbol{+\Delta p_v}
\end{flalign}

Vereinfachung mit $p_{\text{geodätisch}} = const.$ und $p_{\text{dynamisch}} = const.$:
\begin{flalign}
p_1&= p_2+\Delta p_v
\end{flalign}

Kontinuitätsgleichung:
\begin{flalign}
	\dot{V}	&= v*A\\
	\dot{V_1}&= \dot{V_2}\\
	v_1*A_1	&= v_2*A_2
\end{flalign}

Druckverluste:
\begin{flalign}
	\Delta p_v	&= \left(\lambda *\frac{l}{d}+\sum\zeta_i\right)*\frac{\rho}{2}*v^2
\end{flalign}

\textsc{Reynold}szahl:
\begin{flalign}
	Re	&= \frac{v*d*\rho}{\eta} = \frac{v*d}{\nu}
\end{flalign}

K$_V$-Wert:
\begin{flalign}
K_V	&= \dot{V}*\sqrt{\frac{\SI{1}{\bar}}{\Delta p}*\frac{\rho}{\rho_0}}
\end{flalign}

\section{Geräte und Chemikalien}
\label{sec:geraete}

\textcolor{red}{MInipage}

\textbf{Geräte:}
\begin{itemize}
	\begin{minipage}{0.5\textwidth}
		\item Magnetrührer mit Rührfisch
		\item Bechergläser (\SI{100}{\milli \liter}, \SI{400}{\milli \liter})
		\item Messzylinder (\SI{10}{\milli \liter}, \SI{5}{\milli \liter})
		\item Büchnertrichter
		\item Saugflasche
		\item Waschflasche
		\item Filterpapier
		\item Heizpilz
		\item Glasstab
		\item Pipetten
		\item Schutzhandschuhe
	\end{minipage}
	\begin{minipage}{0.5\textwidth}
		\item Filterpapier
		\item Ultraschallbad
		\item Löffel und Spatel
		\item Rundkolben
		\item Rückflusskühler
		\item Wasser- und Vakuumschläuche
		\item Hebebühne
		\item DC-Platte
		\item UV-Licht
		\item Kieselgel (\textsc{Merck} 60)
	\end{minipage}
\end{itemize}

\newpage

\textbf{Proben/Chemikalien:}
\begin{itemize}
	\begin{minipage}{0.55 \textwidth}
		\item \textit{Ansatz:} \SI{1,3}{\gram} bzw. \SI{9}{\milli \mol} 2-Naphthol
		\item destilliertes Wasser
		\item \SI{18}{\milli \liter} \SI{2,5}{\percent}ige Natriumcarbonat-Lösung
		\item \SI{1,7}{\gram} Sulfanilsäure
		\item Eis nach Bedarf
		\item \SI{0,8}{\gram} Natriumnitrit
	\end{minipage}
	\begin{minipage}{0.45 \textwidth}
		\item \SI{2}{\milli \liter} konzentrierte Salzsäure (genutzt wurde zum Teil 1N-Salzsäure)
		\item \SI{8}{\milli \liter} \SI{10}{\percent}ige Natronlauge
		\item \SI{5}{\gram} Natriumchlorid
		\item \SI{10}{\milli \liter} 4:1 Wasser-Ethanol-Mischung
		\item Methanol
		\item 4:1 Essigsäureethylester-Methanol-Mischung
	\end{minipage}
\end{itemize}

\vspace*{5mm}

\textbf{Geräte/Proben/Chemikalien, die im Versuch aus Zeitgründen nicht genutzt wurden:}
\begin{itemize}
	\item Stoff zum Einfärben (Färbegut)
	\item \SI{3}{\percent} von  Färbegut an konzentrierter Schwefelsäure
	\item 20-40 mal die Masse von Färbegut an Wasser
	\item \SI{15}{\percent} von  Färbegut an Natriumsulfat
	\item \SI{30}{\percent}ige Natronlauge
	\item \SI{8}{\milli \liter} \SI{30}{\percent}ige Wasserstoffperoxid
	\item \SI{3}{\gram} Aktivkohle
\end{itemize}





\section{Versuchsdurchführung}
\label{sec:durchfuerung}

\begin{figure}[h!]
	\centering
	\includegraphics[width=0.8\textwidth]{img/versuchsaufbau}
	\caption{Schematischer Versuchsaufbau der Nitrierung}
	\label{fig:versuchsaufbau}
\end{figure}
\FloatBarrier
%Ende
\vspace*{-5mm}

\subsection*{Durchführung:}
Zu Beginn des Versuches wurden \SI{6}{\milli \liter} konzentrierte Schwefelsäure im Rundkolben vorgelegt und mittels Eisbad auf \SI{10}{\celsius} gekühlt. Es folgte die Zugabe der konzentrierten Salpetersäure. Da das Küken des Feststofftrichters zuvor nicht ausreichend auf Funktionalität und benötigtes Schlifffett geprüft wurde, erfolgte zu Beginn keine tropfenweise Zugabe der Salpetersäure. Die Temperatur stieg somit kurzzeitig auf \SI{20}{\celsius} an. Nachdem der Inhalt des Kolbens wieder auf \SI{10}{\celsius} sank, wurde nun tropfenweise weitere Salpetersäure zugegeben. Die Lösung innerhalb des Kolbens färbte sich mit der Zeit gelblich. \\
Die Nitriersäure im Kolben wurde daraufhin auf \SI{15}{\celsius} erwärmt und es wurde mit der Zugabe von Brombenzol in die Lösung begonnen. Diese erfolgte in ca. \SI{2,3}{\minute}-Intervallen à \SI{0,5}{\milli \liter} Brombenzol-Lösung. Ab einer zugegebenen Menge von \SI{2}{\milli \liter} Brombenzol ließ sich im Kolben eine gelblich, weiße Ausflockung beobachten. Nachdem insgesamt \SI{6,5}{\milli \liter} Brombenzol zugegeben wurden, waren vorwiegend um den Rührfisch im Kolben die zuvor beschriebenen Ausflockungen zu erkennen. Während der Zugabe das Brombenzols blieb die Temperatur innerhalb des Kolbens zwischen 26-\SI{35}{\celsius}.\\
Die Lösung wurde daraufhin auf \SI{55}{\celsius}  mittels Wasserbad erwärmt und \SI{45}{\minute} lang auf dieser Temperatur gehalten. Es erfolgte danach ein Abkühlen des Kolbeninhaltes auf Raumtemperatur. Der Kolbeninhalt wurde daraufhin in ein Becherglas mit \SI{60}{\milli \liter} Eiswasser gegeben und Reste innerhalb des Kolbens mit einem Spatel ausgekratzt.\\
\newpage
Anschließend wurde der Inhalt des Becherglases mit einem \textsc{Büchner}-Trichter abgesaugt und es blieb ein gelber Filterkuchen und ein farbloses Filtrat über. Infolgedessen wurde der Filterkuchen in Wasser gegeben und erneut abgesaugt. Das wurde ein weiteres Mal wiederholt und zum Schluss noch einmal abgesaugt. Nach dem dritten Mal absaugen wurde der Filterkuchen auf einem Tonteller getrocknet und anschließend erneut in einen Rundkolben gegeben und mit ca. \SI{30}{\milli \liter}  \SI{95}{\percent}-igen Ethanol umkristallisiert. Der Heizpilz wurde hierbei auf Stufe 3 gestellt.\\
Das aus dem Ethanol ausgefallene 4-Nitrobrombenzol wurde danach heiß in ein Becherglas gegeben und infolgedessen erneut abgesaugt und mit kaltem Ethanol gewaschen. Das Ethanol wurde dafür kurzzeitig in ein Eisbad gestellt. Das Filtrat dieses Absaugens wurde danach in einem Rotationsverdampfer abdestilliert.

\subsection*{Isolierung und Reinigung}
Während dem Umkristallisieren des Rohproduktes wurden ein Teil des Rohproduktes aus 2- und 4- Nitrobrombenzol, sowie Reinstproben von 2-Nitrobrombenzol und 4-Nitrobrombenzol mit \SI{2,7}{\milli \gram}, \SI{4,0}{\milli \gram} und \SI{4,2}{\milli \gram} in jeweils drei Reagenzgläser eingewogen. Jedem Reagenzglas wurde danach jeweils \SI{1}{\milli \liter} Essigsäureethylester zugegeben.\\
Nachdem die Proben für die Dünnschichtchromatografie vorbereitet waren, wurden zwei Fließmittel à \SI{10}{\milli \liter} hergestellt. Diese enthielten ein Gemisch aus Cyclohexan und Essigsäureethylester. Im ersten Fließmittel wurde ein Verhältnis von 8:1 im zweiten ein Verhältnis von 14:1 gewählt. Mittels Kapillaren wurden nun zwei DC-Platten jeweils mit Rohprodukt, 2-Nitrobrombenzol und 4-Nitrobrombenzol betüpfelt. Die Entwicklung der DC-Platten erfolgte in zwei separaten Kammern mit den jeweiligen Fließmitteln. \\
Zum Schluss wurden die entwickelten DC-Platten unter einer UV-Lampe untersucht.\\
Zu einem separaten Termin wurden die Massen der Produkte, sowie deren Schmelzpunkte bestimmt.

\subsection*{Entsorgung:}
Das wässrige Filtrat vom Rohprodukt, sowie die wässrige Phase beim Waschen wurde im Sammelbehälter für halogenhaltige Abfälle entsorgt. Die Überreste der Reagenzgläser mit den Laufmitteln wurden ebenfalls im halogenhaltigen Abfall entsorgt. Ansonsten wurden jegliche Glasapparaturen bzw. Glasgeräte mit Aceton gespült.

\newpage

\section{Ergebnisse}
\label{sec:ergebnisse}

\subsection*{Ausbeute}
Die zu Beginn des Versuches abgemessenen \SI{2,5}{\gram} 3-Nitrophthalsäure und \SI{5}{\milli\liter} 8\%ige Hydrazinhydrat-Lösung entsprachen Stoffmengen von jeweils \SI{12}{\milli \mol}. Die Masse an synthetisierten Luminol wurde trocken mit \anmerkung{\SI{1,88}{\gram}} eingewogen und eine molare Masse von \SI{177}{\gram \per \mole} bestimmt. Aus diesen Angaben wurden die Ausbeute der Luminol-Synthese in Gleichung \ref{gl:ausbeute} bestimmt.\\
\begin{flalign}
	\label{gl:ausbeute}
	\eta 	&= \frac{n_{\text{Luminol}}}{n_{\text{3-Nitrophthalsäure}}} = \frac{m_{\text{Luminol}}}{n_{\text{3-Nitrophthalsäure}}*M_{\text{Luminol}}}\\[2mm]
	&=	\frac{\SI{1,88}{\gram}}{\SI{12}{\milli \mol}*\SI{177}{\gram \per \mole}}\\
	&=\underline{\SI{88,5}{\percent}}
\end{flalign}
Der Wert für die Ausbeute wird unter Abschnitt \ref{sec:diskussion} diskutiert.

\subsection*{Schmelzpunktbestimmung}
Über den Schmelzpunkt ließ sich bestimmen, wie rein das Syntheseprodukt vorliegt. Für das synthetisierte Luminol wurde eine Schmelztemperatur von rund \SI{460}{\celsius} bestimmt. Während der Untersuchung färbte sich das Produkt dabei braun bis schwarz, als würde es verbrennen. Im Vergleich dazu gab eine Literaturrecherche für reines Luminol eine Schmelztemperatur von 329-\SI{332}{\celsius} bzw. 319-\SI{320}{\celsius} an \cite{Luminol_Rompp}. Auffallend während dieser Untersuchung ist die starke Verfärbung des Produktes und ein weit abweichender Schmelzpunkt mit über \SI{100}{\kelvin} unterschied.


\subsection*{Beobachtung der Chemolumineszenz}
Durch Mischung der peroxidischen Kaliumhexacyanoferrat(III)-Lösung mit der Luminol-Lösung entstand eine hellblau leuchtende Lösung.




\newpage
\section{Diskussion der Ergebnisse}
\label{sec:diskussion}

Eine Ausbeute von \SI{15,48}{\percent} erscheint realistisch, da Nelken die ätherischen Öle lediglich als flüchtige Komponenten und nicht Strukturbausteine der Blüten enthalten.\\


Anhand des Vergleiches der $R_f$-Werte zwischen Referenzsubstanz und Produkt lässt sich grob die Reinheit des Produktes bewerten. Da die $R_f$-Werte von 2-Naphthol und Orange II mit 0,918 und 0,411 sich relativ unterschiedlich sind, ist davon auszugehen, dass sich das 2-Naphthol in der Synthese zum Großteil umgesetzt hat. Da jedoch das 2-Naphthol keinen eindeutigen Endpunkt für die Dünnschichtchromatografie aufwies, sollte die genutzte 4:1 Essigsäureethylester-Methanol-Lösung als Laufmittel hinterfragt werden. Für genauere Untersuchungen ist unter Umständen ein alternatives Laufmittel zu verwenden.


%\section{Zusammenfassung und Fazit}
\label{sec:zusammenfassung}


%\newpage
\section{Fehlerbetrachtung}
\label{sec:fehlec|}

Es ist davon auszugehen, dass die größten Fehlerquellen beim Einstellen der Heizleistung des Heizlüfters, des Kühlwasservolumenstromes und dem zeitlich passenden Ablesen der Temperaturen im stationären Zustand des Systems "`Kühlturm"' auftraten. Da sich die Leistung des Heizlüfter nur stufenlos einstellen ließ und die Änderungsrate zwischen der Heizlüfter- Einstellung und der Lufttemperatur sehr war, ließen sich Temperaturen der Luft nur schlecht einstellen. Daher kommt es auch zu den Überschreitungen der Diagrammgrenzen der hx-Diagramme. Weiterhin ließ sich der Kühlwasserstrom weder im Volumen noch in der Temperatur automatisch regeln. Es erforderte daher Feingefühl beim Einstellen des Wasserstromes zusammen mit dem konstant halten der Wassertemperatur. Dieser Umstand könnte durch ein Thermostat behoben werden.\\
 Weiterhin war durch die wechselnden Temperaturen am Computermonitor durch die zuvor genannten Gegebenheiten ein Abschätzen des stationären Zustandes Systems relativ schwierig. Eventuell könnte dieser Punkt behoben werden indem längere Abtastraten der Sensoren dafür sorgen, dass Messwerte in längeren, zeitlichen Abständen aufgenommen werden oder in der Computer Mittelwerte über mehrere Messreihe für bestimmte Zeitintervalle darstellt.

Als weitere Fehlerquelle ist die Dichtigkeit des Kühlturm zu nennen. Dieser zeigte während des Versuches eine geringe Leckage des gekühlten, simulierten Kühlwassers auf und kann somit zum Teil die Messergebnisse verfälschen.

Die Sensoren selbst funktionierten geringfügig Zeitverzögert, zeigten ansonsten für den genutzten Temperaturbereich jedoch keine auffälligen Schwächen auf.

Grundsätzlich liegen die größten Fehlerquellen im Einstellen, der unterschiedlichen Betriebsparameter und nicht in der Messung dieser. Durch prozesstechnische Regelstrukturen könnten diese problematischen Bedingungen eingedämmt werden.

%%Praktikumsskript, Modul ………, Versuch …….., Prof. Musterprof. 
%%DIN 12345, Jahr der Veröffentlichung 
%%Link der Internetseite, Zugriffsdatum 
%%Buchtitel, Autor, Verlag, Veröffentlichungsjahr 
%
%%Literaturverzeichnis Bücher
%\bibliography{Literatur}
%\bibliographystyle{unsrtdin}
%\addcontentsline{toc}{section}{Literaturverzeichnis}



%\section*{Anhang}
\addcontentsline{toc}{section}{Anhang}
\label{sec:anhang}

%Betriebsanweisung

%\chapter*{Eidesstattliche Erklärung}
\label{erklaerung}
Hiermit versichere ich, die vorliegende Seminararbeit selbstständig und nur unter Verwendung der von mir angegebenen Quellen und Hilfsmittel verfasst zu haben. Sowohl inhaltlich als auch wörtlich entnommene Inhalte wurden als solche kenntlich gemacht. Die Arbeit hat in dieser oder vergleichbarer Form noch keinem anderem Prüfungsgremium vorgelegen. \\
\\[1.5cm]
Datum:	\hrulefill\enspace Unterschrift: \hrulefill
\\[3.5cm]
\addcontentsline{toc}{chapter}{Selbstständigkeitserklärung}

\end{document}
